%%%%%
%
%% Dans ce fichier, les �ventuels accents sont cod�s
%% en ISO-8859-15 (aka latin9)
%% La ligne suivante donne trois exemples 
%% e accent aigu : �, e accent grave : �, e-dans-l'o : �
\section{Indicateur de qualit� du signal lidar}

\itemstyle{Op\'erateur~:} \operfmt{Qualite\_profil\_lidar}
\index{Qualite\_profil\_lidar}

%%%%%OBJET%%%%
\itemstyle{Objet~:}
Calcul d'un indicateur de qualit� du signal lidar par rapport � la r�tro-diffusion mol�culaire.

%%%%%ARGUMENTS%%%%
\bigskip
%% Dans ce fichier, les �ventuels accents sont cod�s
%% en ISO-8859-15 (aka latin9)
%% La ligne suivante donne trois exemples 
%% e accent aigu : �, e accent grave : �, e-dans-l'o : �
%%%%%ENTREES%%%%
\begin{tabarg}
\multicolumn{3}{l}{\titretab {Entr\'ees}}\\
\hline
$X$ & Vecteur, \texttt{float[$N_{level}$]} & Signal lidar \\
$X_{mol}$ & Vecteur, \texttt{float[$N_{level}$]} & Signal de r�tro-diffusion mol�culaire \\
$Z$ & Vecteur, \texttt{float[$N_{level}$]} & Altitudes r�elles des portes du lidar (m)\\
$Z_{ref}$ & Vecteur, \texttt{float} & Altitude de r�f�rence du signal mol�culaire  (m)\\
\end{tabarg}

\medskip
%%%%%SORTIES%%%%
\begin{tabarg}
\multicolumn{3}{l}{\titretab {Sorties}}\\
\hline
$Q$ & Vecteur, \texttt{float} & Indicateur de qualit� du profil lidar \\
\end{tabarg}



%%%%%FORMULE%%%%
\itemstyle{Formule}

%% Formule de l'operateur : Qualite_profil_lidar
%% Dans ce fichier, les �ventuels accents sont cod�s
%% en ISO-8859-15 (aka latin9)
%% La ligne suivante donne trois exemples 
%% e accent aigu : �, e accent grave : �, e-dans-l'o : �

Cet op�rateur calcule la pente du signal lidar $X$ dans une tranche d'altitude de 1500 m d'�paisseur autour de l'altitude $Z_{ref}$. Le rapport de cette pente et de celle du signal mol�culaire $X_{mol}$ repr�sente la qualit� du signal $Q$.



%%%%%AUTEUR%%%%
\itemstyle{Auteur~:}
CNRM/GMEI/TRAMM

%%%%%REFERENCE%%%%
\itemstyle{R\'ef\'erences~:}
Utilis� dans la fili�re BASIC du LOA.
A. Mortier 2013 \cite{THESE-MORTIER-2013}.
